% =============================================================
% IEEE iSPEC 2026 Conference Paper
% Title: A Three-Stage Deep Learning Pipeline for Automated
%        Cell-Level Extraction, Defect Triage, and Fine-Grained
%        Segmentation in Photovoltaic Modules
% Template: IEEE Conference Template (062824)
% =============================================================

\documentclass[conference]{IEEEtran}
\IEEEoverridecommandlockouts
% The preceding line is only needed to identify funding in the
% first footnote. If that is unneeded, please comment it out.

\usepackage{cite}
\usepackage{amsmath,amssymb,amsfonts}
\usepackage{algorithmic}
\usepackage{graphicx}
\usepackage{textcomp}
\usepackage{xcolor}
\usepackage{booktabs}   % for professional tables
\usepackage{multirow}   % for multi-row table cells
\usepackage{url}        % for URLs in references

\def\BibTeX{{\rm B\kern-.05em{\sc i\kern-.025em b}\kern-.08em
    T\kern-.1667em\lower.7ex\hbox{E}\kern-.125emX}}

% ---- Convenience command for TBD placeholders ----
\newcommand{\tbd}[1][\,]{\textbf{[TBD#1]}}

\begin{document}

% ================================================================
%  TITLE
% ================================================================
\title{A Three-Stage Deep Learning Pipeline for Automated\\
Cell-Level Extraction, Defect Triage, and Fine-Grained\\
Segmentation in Photovoltaic Modules%
\thanks{This work was carried out at the Indian Institute of
Technology Mandi. The authors gratefully acknowledge the
support of IIT Mandi and the guidance of their research
supervisor.}
}

% ================================================================
%  AUTHORS  -- fill in actual names, affiliations and emails
% ================================================================
\author{%
\IEEEauthorblockN{1\textsuperscript{st} \textbf{[Author 1 Name]}}
\IEEEauthorblockA{\textit{\textbf{[Department]}} \\
\textit{Indian Institute of Technology Mandi}\\
Mandi, Himachal Pradesh, India \\
\textbf{[email@iitmandi.ac.in]}}
\and
\IEEEauthorblockN{2\textsuperscript{nd} \textbf{[Author 2 Name]}}
\IEEEauthorblockA{\textit{\textbf{[Department]}} \\
\textit{Indian Institute of Technology Mandi}\\
Mandi, Himachal Pradesh, India \\
\textbf{[email@iitmandi.ac.in]}}
}

\maketitle

% ================================================================
%  ABSTRACT
% ================================================================
\begin{abstract}
Automated visual inspection of photovoltaic (PV) modules is
essential for large-scale solar farm quality control and
predictive maintenance, yet existing approaches are typically
evaluated only on pre-cropped, single-cell images and collapse
fine-grained defect taxonomies into coarse binary labels.
This paper proposes a two-stage pipeline (with a roadmap to
a third segmentation stage) that operates directly on
module-level electroluminescence (EL) imagery---including
full modules and partial or occluded captures---spanning both
monocrystalline and polycrystalline cell types.
Stage~1 performs cell extraction from raw module images
using YOLOv8-OBB (oriented bounding boxes), adopted after a
systematic evaluation of eleven classical OpenCV-based
approaches across two dataset tracks
(CellExtractOpenCV and FullModuleOpenCV), none of which
generalised robustly across module formats and imaging
conditions.
Stage~2 performs binary quality triage (OK vs.\ B-grade)
using CNN backbones (ResNet, EfficientNet) augmented with a
Convolutional Block Attention Module (CBAM), with
hyperparameters tuned via Optuna.
A third stage for pixel-level defect segmentation---using a
principled 29$\to$4 defect-class consolidation and a
comparison of convolutional (U-Net family) against
transformer-based (SegFormer) architectures---is under active
development.
We report \tbd{} mAP for Stage~1 cell extraction and
\tbd{} AUROC for Stage~2 triage.
\end{abstract}

\begin{IEEEkeywords}
photovoltaic inspection, electroluminescence imaging,
YOLOv8-OBB, oriented bounding boxes, defect classification,
CBAM, attention mechanism, Optuna, deep learning, solar energy
\end{IEEEkeywords}

% ================================================================
%  SECTION I --- INTRODUCTION
% ================================================================
\section{Introduction}

The global PV industry is undergoing rapid capacity expansion,
with cumulative installed solar capacity projected to exceed
10~TW by~2030. Maintaining cell-level quality throughout the
module manufacturing lifecycle and in-field operation is a
critical bottleneck: manual electroluminescence (EL) imaging
inspection is labour-intensive, subjective, and fails to scale
to the throughput requirements of modern manufacturing lines
or drone-based fleet inspection.

Three unresolved gaps motivate this work.
\textit{First}, most published EL defect classification and
segmentation studies---including benchmark work on the ELPV
dataset~\cite{deitsch2019}---assume that individual cell crops
are provided as input.
In realistic deployment, raw module-level images arrive with
arbitrary orientation, partial occlusion, and perspective
distortion; no robust plug-in pipeline exists to bridge this gap.
\textit{Second}, binary ``defective/non-defective''
classification is the prevailing paradigm.
While sufficient for pass/fail sorting, it provides no
information about \textit{which} failure mechanism is present
or \textit{how severe} it is---information needed for
warranty decisions, rework prioritisation, and predictive
degradation modelling.
\textit{Third}, fine-grained EL defect ontologies enumerate
20--30 subtypes; yet many subtypes are rare and visually
ambiguous, making direct multi-class segmentation
statistically intractable without a principled class
consolidation strategy.

The contributions of this paper are:
\begin{enumerate}
  \item An end-to-end pipeline operating on \textbf{raw
        module-level EL images} (full or partial, mono- and
        polycrystalline), including perspective correction and
        cell normalisation.
  \item A rigorous ablation of \textbf{fifteen
        OpenCV-based classical extraction iterations} (eleven
        on partial/single-cell images, four on full-module
        images), establishing root-cause failure modes and
        motivating the adoption of deep-learning detection.
  \item A \textbf{YOLOv8-OBB} detector for oriented cell
        extraction, trained on 2{,}212 module images
        (1{,}769 train / 443 val), serving as the unified
        Stage~1 detector for both full-module and
        partial-module inputs.
  \item A CBAM-augmented, Optuna-tuned \textbf{binary triage
        stage} benchmarked across ResNet and EfficientNet
        backbones.
  \item A forward-looking description of the Stage~3
        segmentation roadmap: a defect-mechanism-informed
        \textbf{29$\to$4 class consolidation} and a planned
        comparison of CNN vs.\ transformer segmentation
        architectures.
\end{enumerate}

% ================================================================
%  SECTION II --- RELATED WORK
% ================================================================
\section{Related Work}

\subsection{PV Defect Datasets}

The ELPV dataset~\cite{deitsch2019} is the canonical benchmark
for EL-based defect classification: 2{,}624 grayscale
cell-level images (300$\times$300~px), labelled with defect
probability and cell technology (mono/poly). Its widespread
adoption has shaped the research paradigm toward pre-cropped
single-cell inputs, leaving the module-to-cell extraction
problem largely unaddressed.
For Stage~3, a multi-class pixel-annotated EL dataset
encompassing 29 defect subtypes with polygon-level
annotations is used~\cite{dataset_tbd}.

\subsection{Classical Cell Extraction from Module Images}

Classical approaches to PV cell grid detection exploit the
periodic spatial signal of busbars and inter-cell boundaries.
Standard methods include Hough-transform-based line
detection~\cite{hough_lines}, morphological line
extraction~\cite{morph_lines}, projection
profiling~\cite{projection_profiles}, and
DBSCAN-assisted line clustering.
While elegant, these approaches assume approximately
axis-aligned images, low sensor noise, and consistent
inter-cell contrast---conditions that fail for tilted,
partially occluded, or field-degraded module captures.

\subsection{Deep Learning for Cell Detection}

YOLO-family models~\cite{redmon2016,jocher2023} demonstrate
state-of-the-art throughput--accuracy trade-offs in
industrial inspection.
The YOLOv8-OBB variant extends axis-aligned detection with
oriented bounding boxes, directly predicting a rotation angle
alongside standard box parameters---making it suited to
arbitrary-orientation PV cell extraction without
pre-rectification.

\subsection{Defect Classification and Segmentation}

ResNet~\cite{he2016} and EfficientNet~\cite{tan2019} are the
dominant backbones for EL defect classification.
CBAM~\cite{woo2018} augments any CNN backbone with lightweight
channel-wise and spatial-wise attention.
For segmentation, U-Net~\cite{ronneberger2015} and
U-Net++~\cite{zhou2018} dominate pixel-level inspection.
Transformer-based SegFormer~\cite{xie2021} offers superior
global context modelling, advantageous for spatially diffuse
defects such as inactive regions and corrosion halos.

% ================================================================
%  SECTION III --- DATASET
% ================================================================
\section{Dataset}

\subsection{Stage 1 --- Module-Level EL Images}

The Stage~1 dataset consists of high-resolution JPEG EL images
of commercial PV modules acquired under controlled imaging
conditions.
The dataset spans two module formats: 144-cell (6-string)
and 208-cell (8-string) configurations.
\begin{itemize}
  \item \textbf{Total images:} 2{,}212 full-module EL captures.
  \item \textbf{Split:} 80\,\% training (1{,}769~images) /
        20\,\% validation (443~images); split at the module
        level to prevent cell-level data leakage.
  \item \textbf{Classical CV subsets:} ARTS, SDLE, and BMRK
        datasets were used during the classical OpenCV
        exploration phase, each representing different imaging
        protocols and module manufacturers.
\end{itemize}

\subsection{Stage 2 --- Cell-Level Triage Images}

Cells for binary triage are drawn from the ELPV
dataset~\cite{deitsch2019} combined with cell crops extracted
by the Stage~1 pipeline.
Labels are binarised: OK (functional) vs.\ B-grade (defective).
A module-level train/val/test split is applied to avoid
spatial correlation across folds.

\subsection{Stage 3 --- Pixel-Annotated B-Grade Cells}

A multi-class EL segmentation dataset with polygon-level
annotations covering 29 defect subtypes~\cite{dataset_tbd}
provides training data for Stage~3.
Table~\ref{tab:taxonomy} summarises the 29$\to$4 superclass
consolidation motivated by physical failure mechanisms.

\begin{table}[t]
  \caption{Defect taxonomy: 29$\to$4 superclass consolidation}
  \label{tab:taxonomy}
  \begin{center}
  \renewcommand{\arraystretch}{1.15}
  \begin{tabular}{lp{2.8cm}p{1.5cm}}
    \toprule
    \textbf{Superclass} & \textbf{Physical Mechanism} &
    \textbf{Class IDs} \\
    \midrule
    1 --- Crack &
      Mechanical fracture of Si wafer &
      10, 14, 28 \\
    2 --- Inactive Area &
      Electrically disconnected region &
      11, 17, 20 \\
    3 --- Corrosion/Metal. &
      Finger/busbar degradation, oxidation &
      15, 16, 18, 26 \\
    4 --- Material/Surface &
      Contamination, discoloration &
      12, 13, 19, 25, 27 \\
    \bottomrule
  \end{tabular}
  \end{center}
\end{table}

\subsection{General Preprocessing}

All EL images are loaded as grayscale (single-channel) 8-bit
images. The common preprocessing chain includes:
(i)~median denoising (3$\times$3 kernel) to suppress
salt-and-pepper noise;
(ii)~CLAHE enhancement (clipLimit~=~2.0,
tileGridSize~=~8$\times$8) to normalise non-uniform EL emission;
(iii)~perspective correction via \texttt{cv2.warpPerspective}
to warp each extracted cell to a canonical upright view; and
(iv)~standard normalisation (mean subtraction, std division)
before deep learning inference.

% ================================================================
%  SECTION IV --- METHODOLOGY
% ================================================================
\section{Methodology}

\subsection{Stage 1 --- Cell Extraction}
\label{sec:stage1}

\subsubsection{Phase A: Classical OpenCV (CellExtractOpenCV)}

Eleven notebooks systematically explored classical extraction
on partial-module and single-cell images
(preprocessing1--preprocessing11).
Table~\ref{tab:classical_failure} summarises each approach and
its key limitation.

\begin{table}[t]
  \caption{Classical OpenCV extraction: failure summary}
  \label{tab:classical_failure}
  \begin{center}
  \renewcommand{\arraystretch}{1.1}
  \small
  \begin{tabular}{lp{1.9cm}p{2.4cm}}
    \toprule
    \textbf{Notebook} & \textbf{Technique} &
    \textbf{Key Limitation} \\
    \midrule
    pp1  & CLAHE + intensity profiling & Exploratory; no extraction \\
    pp2  & Morph.\ grid + persp.\ warp & Single image; not generalised \\
    pp3  & Inversion + contour & Conflates cell with bright defects \\
    pp4  & Batch \texttt{extract\_middle\_cell} & 5--10\,\% skip rate \\
    pp5  & --                       & YOLO pivot notebook \\
    pp6  & Gaussian + Canny         & Fragmented contours \\
    pp7  & NL-Means + Hough         & Image-specific thresholds \\
    pp8  & Same on multi-cell       & Requires a priori line count \\
    pp9  & Otsu + Canny             & Format-specific count limit \\
    pp10 & Otsu + contour polygon   & Fails on multi-cell images \\
    pp11 & Hough + DBSCAN           & Dataset-specific $\epsilon$ tuning \\
    \bottomrule
  \end{tabular}
  \end{center}
\end{table}

Four systematic failure modes were identified:
\begin{enumerate}
  \item \textit{Hyperparameter brittleness} --- CLAHE clip
        limits, Hough thresholds, and DBSCAN~$\epsilon$ values
        tuned for one image break on the next.
  \item \textit{Module format sensitivity} --- the projection
        peak detector counted 9 horizontal lines for one module
        and 8 for the next (same physical format), yielding
        208 vs.\ 182 cells.
  \item \textit{Contrast variability} --- low-contrast
        degraded regions cause grid detectors to miss busbar
        lines below the adaptive threshold.
  \item \textit{Perspective intolerance} --- morphological
        kernels assume axis-aligned structure; even
        5--10$^\circ$ module tilt causes systematic
        misdetection.
\end{enumerate}

\subsubsection{Phase B: Classical OpenCV (FullModuleOpenCV)}

A separate four-notebook track (fullmodulev1--fullmodulev4)
targeted full-resolution module images.
\texttt{fullmodulev2} achieved the first successful full-module
extraction (144~cells from one module) using projection-profile
peak detection after automatic dark-edge cropping.
\texttt{fullmodulev4} processed 10~modules with DBSCAN-based
line merging but still required manual cell-count verification
per module and provided no perspective-correction capability.

\subsubsection{Phase C: YOLOv8-OBB (Final Approach)}

YOLOv8-OBB (Ultralytics~8.3.235)~\cite{jocher2023} was
adopted as the unified Stage~1 detector.
Each detection outputs five parameters:
$(c_x, c_y, w, h, \theta)$, where $\theta$ is the rotation
angle, enabling extraction of arbitrarily oriented cells
without pre-rectification.
Key advantages over classical methods: (i)~format agnosticism
through learned features; (ii)~rotation robustness via
oriented bounding boxes; (iii)~scale invariance from the
multi-scale feature pyramid; (iv)~unified coverage of
full-module, partial-module, and single-cell inputs with one
model.

The post-processing pipeline converts OBB outputs to
rectified cell crops:
\begin{equation}
  (c_x, c_y, w, h, \theta)
  \;\xrightarrow{\text{corner extraction}}\;
  M_{\text{persp}}
  \;\xrightarrow{\texttt{warpPerspective}}\;
  \hat{I}_{\text{cell}}
  \label{eq:postproc}
\end{equation}
where $\hat{I}_{\text{cell}}$ is the canonical
perspective-corrected crop.

Fig.~\ref{fig:yolo_detections} shows representative YOLOv8-OBB
detections on a full-module EL image and resulting
perspective-corrected cell crops.

\begin{figure}[t]
  \centering
  \includegraphics[width=\columnwidth]{fig_yolo_detections.png}
  \caption{YOLOv8-OBB detections overlaid on a full-module
           EL image (top) and resulting perspective-corrected
           cell crops (bottom). Oriented bounding boxes
           accommodate module tilt.}
  \label{fig:yolo_detections}
\end{figure}

\subsection{Stage 2 --- Binary Quality Triage}
\label{sec:stage2}

\subsubsection{Architecture}

Extracted and rectified cells from Stage~1 are classified as
structurally intact (OK) or degraded (B-grade) using a
pretrained CNN backbone (ResNet or EfficientNet) augmented with
a CBAM block~\cite{woo2018} appended after the final
convolutional stage, followed by global average pooling, a
dropout-regularised fully connected layer, and a sigmoid
output.

CBAM applies dual attention sequentially.
\textit{Channel attention} reweights feature channels:
\begin{equation}
  \mathbf{M}_c(\mathbf{F}) = \sigma\!\left(
    W_1\!\left(\text{ReLU}\!\left(W_0\,\text{GAP}(\mathbf{F})\right)\right) +
    W_1\!\left(\text{ReLU}\!\left(W_0\,\text{GMP}(\mathbf{F})\right)\right)
  \right)
  \label{eq:cbam_channel}
\end{equation}
\textit{Spatial attention} then localises discriminative regions:
\begin{equation}
  \mathbf{M}_s(\mathbf{F}') = \sigma\!\left(
    f^{7\times7}\!\!\left(
      \left[\text{AvgPool}(\mathbf{F}');\,\text{MaxPool}(\mathbf{F}')\right]
    \right)\right)
  \label{eq:cbam_spatial}
\end{equation}
where $\sigma$ is sigmoid, $W_0, W_1$ are MLP weights with
reduction ratio~$r$, $f^{7\times7}$ is a $7{\times}7$
convolution, and $\mathbf{F}' = \mathbf{M}_c \otimes \mathbf{F}$.
Since defect features in EL images occupy typically $<$5\,\%
of cell area, CBAM's dual attention suppresses irrelevant
background without requiring localisation annotations.

\subsubsection{Hyperparameter Optimisation via Optuna}

Table~\ref{tab:optuna} lists the hyperparameter search space.
The optimisation objective is validation AUROC, chosen for its
threshold-invariant behaviour under class imbalance.

\begin{table}[t]
  \caption{Optuna hyperparameter search space for Stage~2}
  \label{tab:optuna}
  \begin{center}
  \renewcommand{\arraystretch}{1.1}
  \begin{tabular}{ll}
    \toprule
    \textbf{Hyperparameter} & \textbf{Search Space} \\
    \midrule
    Learning rate & LogUniform$[10^{-5},\;10^{-2}]$ \\
    Backbone & \{ResNet-18, ResNet-50, EffNet-B0, EffNet-B3\} \\
    CBAM reduction ratio $r$ & \{4, 8, 16\} \\
    Dropout rate & Uniform$[0.1,\;0.5]$ \\
    Batch size & \{16, 32, 64\} \\
    Optimizer & \{Adam, AdamW, SGD+momentum\} \\
    Objective & Validation AUROC \\
    \bottomrule
  \end{tabular}
  \end{center}
\end{table}

\subsection{Stage 3 --- Defect Segmentation (Roadmap)}
\label{sec:stage3}

Stage~3 takes B-grade cell crops from Stage~2 as input and
produces pixel-level defect maps over the four consolidated
superclasses (Table~\ref{tab:taxonomy}).
Five architectures are planned for comparison:
U-Net~\cite{ronneberger2015}, U-Net++~\cite{zhou2018},
Attention U-Net, SegFormer-B0, and
SegFormer-B2~\cite{xie2021}.
The loss function will be a combined Dice + Focal loss weighted
by inverse class frequency to address severe class imbalance;
Tversky loss is an additional candidate.
Experiments are ongoing.

\subsection{Pipeline Integration}

Fig.~\ref{fig:pipeline} illustrates the end-to-end cascaded
inference flow.
A known risk is error propagation: missed cells at Stage~1
create false negatives for Stage~2; misaligned crops produce
distorted Stage~3 inputs.
A planned ablation will compare pipeline performance with
ground-truth vs.\ YOLO-predicted cell crops to quantify
upstream error contribution.

\begin{figure}[t]
  \centering
  \includegraphics[width=\columnwidth]{fig_pipeline_overview.png}
  \caption{End-to-end pipeline: raw module EL image
           $\rightarrow$ YOLOv8-OBB extraction
           $\rightarrow$ perspective correction
           $\rightarrow$ CBAM-CNN triage
           $\rightarrow$ (Stage~3, ongoing) segmentation
           $\rightarrow$ per-cell defect map.}
  \label{fig:pipeline}
\end{figure}

% ================================================================
%  SECTION V --- EXPERIMENTAL SETUP
% ================================================================
\section{Experimental Setup}

\textbf{Hardware:} \tbd[: GPU model and VRAM].

\textbf{Software stack:}
Python~3.x, PyTorch~\tbd{}, Ultralytics~8.3.235
(YOLOv8-OBB), \texttt{timm}~\tbd{} (backbone weights),
Optuna~\tbd{}, OpenCV~4.12.0.88, NumPy~1.26.4,
scikit-learn~\tbd{}, Matplotlib~3.8.4.

Table~\ref{tab:training_config} summarises the per-stage
training configuration.

\begin{table}[t]
  \caption{Training configuration per stage}
  \label{tab:training_config}
  \begin{center}
  \renewcommand{\arraystretch}{1.1}
  \small
  \begin{tabular}{lll}
    \toprule
    \textbf{Parameter} &
    \textbf{Stage 1 (YOLO-OBB)} &
    \textbf{Stage 2 (CNN)} \\
    \midrule
    Framework    & Ultralytics     & PyTorch + timm \\
    Input res.   & \tbd{}          & \tbd{} \\
    Epochs       & \tbd{}          & Optuna-selected \\
    Batch size   & \tbd{}          & Optuna-selected \\
    Optimizer    & \tbd{}          & Optuna-selected \\
    LR schedule  & Cosine          & Optuna-selected \\
    Early stop   & patience~\tbd{} & patience~\tbd{} \\
    Random seeds & $\geq 3$        & $\geq 3$ \\
    \bottomrule
  \end{tabular}
  \end{center}
\end{table}

\textbf{Reproducibility:} All experiments run with $\geq$3
independent random seeds. Results are reported as
mean~$\pm$~std across seeds. Data splits are fixed before any
training, reproduced via \texttt{random.seed(<seed>)}.

% ================================================================
%  SECTION VI --- RESULTS
% ================================================================
\section{Results}

\subsection{Stage 1: Cell Extraction}

Table~\ref{tab:stage1_results} compares YOLOv8-OBB against the
best classical baseline (fullmodulev4).

\begin{table}[t]
  \caption{Stage~1 cell extraction performance}
  \label{tab:stage1_results}
  \begin{center}
  \renewcommand{\arraystretch}{1.1}
  \small
  \begin{tabular}{llcccc}
    \toprule
    \textbf{Model} & \textbf{Type} &
    \textbf{mAP@.5} & \textbf{mAP@.5:.95} &
    \textbf{Compl.} & \textbf{O-Seg} \\
    \midrule
    YOLO-OBB & Full-144 & \tbd{} & \tbd{} & \tbd{} & \tbd{} \\
    YOLO-OBB & Full-208 & \tbd{} & \tbd{} & \tbd{} & \tbd{} \\
    YOLO-OBB & Partial  & \tbd{} & \tbd{} & \tbd{} & \tbd{} \\
    Classical (best) & Full & N/A & N/A & $\sim$90\,\%* & High \\
    \bottomrule
    \multicolumn{6}{l}{\footnotesize *Estimated; manual
    verification required per module.} \\
    \multicolumn{6}{l}{\footnotesize Compl.\,=\,Completeness
    rate; O-Seg\,=\,Over-segmentation rate.}
  \end{tabular}
  \end{center}
\end{table}

\subsection{Stage 2: Binary Quality Triage}

Table~\ref{tab:stage2_results} shows the ablation across
backbone and CBAM settings.
Rows~1--2 isolate the CBAM contribution (same backbone, same
LR, with/without CBAM) from the Optuna tuning contribution
in rows~3--5.

\begin{table}[t]
  \caption{Stage~2 binary triage performance}
  \label{tab:stage2_results}
  \begin{center}
  \renewcommand{\arraystretch}{1.1}
  \small
  \begin{tabular}{llcccc}
    \toprule
    \textbf{Backbone} & \textbf{CBAM} &
    \textbf{Optuna} & \textbf{Acc.} &
    \textbf{F1} & \textbf{AUROC} \\
    \midrule
    ResNet-18  & No  & No  & \tbd{} & \tbd{} & \tbd{} \\
    ResNet-18  & Yes & No  & \tbd{} & \tbd{} & \tbd{} \\
    ResNet-50  & Yes & Yes & \tbd{} & \tbd{} & \tbd{} \\
    EffNet-B0  & Yes & Yes & \tbd{} & \tbd{} & \tbd{} \\
    EffNet-B3  & Yes & Yes & \tbd{} & \tbd{} & \tbd{} \\
    \bottomrule
  \end{tabular}
  \end{center}
\end{table}

Results will be reported separately for monocrystalline and
polycrystalline cells, as defect visibility and class
distribution differ by technology.

% ================================================================
%  SECTION VII --- DISCUSSION
% ================================================================
\section{Discussion}

\subsection{Why Classical OpenCV Failed at Scale}

The 15-notebook preprocessing journey reveals four root causes:
(1)~hyperparameter brittleness across imaging conditions;
(2)~module format sensitivity where projection-peak detectors
miscounted grid lines between modules of the same physical
format (208 vs.\ 182 cells);
(3)~contrast variability causing morphological grid detectors
to miss busbar lines in degraded regions; and
(4)~perspective intolerance of axis-aligned morphological
kernels.
YOLOv8-OBB subsumes all four sensitivities into learned model
weights, trading hand-crafted rules for a labelled training
data requirement.

\subsection{Unified Detector Architecture}

A alternative design would use separate detectors for
full-module vs.\ partial-module inputs.
The fullmodulev4 experience demonstrates that even within the
full-module track, format-specific tuning is unavoidable
classically.
A single YOLOv8-OBB model trained on a mix of image types
generalises across both input categories, reducing deployment
complexity and maintenance burden.

\subsection{CBAM for Defect-Aware Triage}

Defect features in EL images are spatially localised and
spectrally discriminative---appearing as dark regions on a
brighter background.
CBAM's channel attention selects defect-relevant feature maps
while spatial attention suppresses background; this is
expected to provide measurable F1 and AUROC gains (rows~1
vs.~2 in Table~\ref{tab:stage2_results}) without requiring
explicit defect localisation annotations during training.

\subsection{Limitations}

\begin{itemize}
  \item \textit{Domain gap}: Training data are lab EL images;
        field-collected EL introduces compression artefacts,
        variable exposure, and ambient light.
  \item \textit{Severity scoring absent}: The current pipeline
        produces a binary triage decision but no continuous
        severity score, limiting actionability for maintenance
        prioritisation.
  \item \textit{Unresolved taxonomy}: Several of the 29
        original defect class IDs remain unmapped; their
        resolution (merge vs.\ drop) affects reported class
        frequencies.
\end{itemize}

% ================================================================
%  SECTION VIII --- CONCLUSION AND FUTURE WORK
% ================================================================
\section{Conclusion and Future Work}

We have presented a two-stage automated inspection pipeline for
PV modules operating on raw module-level EL imagery.
A systematic ablation of 15 classical OpenCV approaches
established that grid-line detection cannot robustly scale to
deployment-level diversity of imaging conditions, module
formats, and orientations, motivating YOLOv8-OBB as a unified
Stage~1 cell extractor.
Stage~2 binary triage uses CBAM-augmented CNNs with
Optuna-based hyperparameter search to address the coarse-label
limitation without localisation supervision.

Planned future work includes:
(1)~completing Stage~3 pixel-level defect segmentation with the
29$\to$4 consolidated taxonomy and CNN vs.\ transformer
comparison;
(2)~severity scoring by mapping segmentation outputs to
power-loss metrics from IV-curve characterisation;
(3)~instance-level defect localisation for counting and spatial
clustering;
(4)~domain adaptation from lab EL to field-acquired imagery;
and
(5)~edge deployment benchmarking on embedded GPU hardware
(e.g., NVIDIA Jetson Orin) for real-time manufacturing
integration.

% ================================================================
%  ACKNOWLEDGMENT
% ================================================================
\section*{Acknowledgment}

This work was supported by the Indian Institute of Technology
Mandi. The authors thank their research supervisor for
invaluable guidance and support throughout this project.

% ================================================================
%  REFERENCES
% ================================================================
\begin{thebibliography}{00}

\bibitem{deitsch2019}
S.\ Deitsch, V.\ Christlein, S.\ Berger, C.\ Buerhop-Lutz,
A.\ Maier, F.\ Gallwitz, and C.\ Riess,
``Automatic classification of defective photovoltaic module
cells in electroluminescence images,''
\textit{Solar Energy}, vol.~185, pp.~455--468, 2019.

\bibitem{he2016}
K.\ He, X.\ Zhang, S.\ Ren, and J.\ Sun,
``Deep residual learning for image recognition,''
in \textit{Proc.\ IEEE Conf.\ Comput.\ Vis.\ Pattern
Recognit.\ (CVPR)}, 2016, pp.~770--778.

\bibitem{tan2019}
M.\ Tan and Q.\ V.\ Le,
``EfficientNet: Rethinking model scaling for convolutional
neural networks,''
in \textit{Proc.\ Int.\ Conf.\ Mach.\ Learn.\ (ICML)},
2019, pp.~6105--6114.

\bibitem{woo2018}
S.\ Woo, J.\ Park, J.-Y.\ Lee, and I.\ S.\ Kweon,
``CBAM: Convolutional Block Attention Module,''
in \textit{Proc.\ Eur.\ Conf.\ Comput.\ Vis.\ (ECCV)},
2018, pp.~3--19.

\bibitem{akiba2019}
T.\ Akiba, S.\ Sano, T.\ Yanase, T.\ Ohta, and M.\ Koyama,
``Optuna: A next-generation hyperparameter optimization
framework,''
in \textit{Proc.\ ACM SIGKDD Int.\ Conf.\ Knowl.\ Disc.\
Data Min.\ (KDD)}, 2019, pp.~2623--2631.

\bibitem{xie2021}
E.\ Xie, W.\ Wang, Z.\ Yu, A.\ Anandkumar, J.\ M.\ Alvarez,
and P.\ Luo,
``SegFormer: Simple and efficient design for semantic
segmentation with transformers,''
in \textit{Adv.\ Neural Inf.\ Process.\ Syst.\ (NeurIPS)},
2021.

\bibitem{ronneberger2015}
O.\ Ronneberger, P.\ Fischer, and T.\ Brox,
``U-Net: Convolutional networks for biomedical image
segmentation,''
in \textit{Proc.\ Int.\ Conf.\ Med.\ Image Comput.\
Comput.-Assist.\ Interv.\ (MICCAI)}, 2015, pp.~234--241.

\bibitem{zhou2018}
Z.\ Zhou, M.\ M.\ Rahman Siddiquee, N.\ Tajbakhsh, and
J.\ Liang,
``UNet++: A nested U-Net architecture for medical image
segmentation,''
in \textit{Proc.\ Deep Learn.\ Med.\ Image Anal.\
(DLMIA/MICCAI)}, 2018.

\bibitem{redmon2016}
J.\ Redmon, S.\ Divvala, R.\ Girshick, and A.\ Farhadi,
``You only look once: Unified, real-time object detection,''
in \textit{Proc.\ IEEE Conf.\ Comput.\ Vis.\ Pattern
Recognit.\ (CVPR)}, 2016, pp.~779--788.

\bibitem{jocher2023}
G.\ Jocher \textit{et al.},
``Ultralytics YOLOv8,'' 2023.
[Online]. Available: \url{https://github.com/ultralytics/ultralytics}

\bibitem{hough_lines}
R.\ O.\ Duda and P.\ E.\ Hart,
``Use of the Hough transformation to detect lines and curves
in pictures,''
\textit{Commun.\ ACM}, vol.~15, no.~1, pp.~11--15, 1972.

\bibitem{morph_lines}
P.\ Soille,
\textit{Morphological Image Analysis: Principles and
Applications}, 2nd~ed.
Berlin: Springer, 2003.

\bibitem{projection_profiles}
K.\ Y.\ Wong,
``Document analysis and recognition using projection
profiles,''
\textit{Pattern Recognit.}, vol.~15, no.~3, pp.~191--199, 1982.

\bibitem{dataset_tbd}
\textbf{[Dataset citation --- to be confirmed upon permission
approval]}, unpublished.

\end{thebibliography}

\end{document}
